\chapter{Fix 13: Vafa-Witten Bounds for Adjoint QCD}
\label{ch:fix13_vafa_witten}

\section{Positivity of the Measure}
A major obstacle in fermionic lattice gauge theories is the "Sign Problem". For the Adjoint Interpolation to define a valid path in the space of theories, the underlying measure must remain positive definite for all masses $M \in [0, \infty)$.

\begin{theorem}[Positivity of Adjoint Determinant]
\label{thm:adjoint_positivity}
Let $D_W(M)$ be the Wilson-Dirac operator in the adjoint representation of the gauge group $G$. For any gauge configuration $U$, the determinant is non-negative:
\begin{equation}
\det(D_W(M)) \ge 0
\end{equation}
\end{theorem}

\begin{proof}
The adjoint representation is real. Therefore, the eigenvalues of the Dirac operator come in complex conjugate pairs $(\lambda, \bar{\lambda})$ or are real.
\begin{enumerate}
    \item The spectrum of $\gamma_5 D_W$ is real because $\gamma_5 D_W$ is Hermitian (provided $M$ is real).
    \item Since $[D_W, C] = 0$ where $C$ is the charge conjugation matrix and the representation is real, the spectrum is doubly degenerate.
    \item Thus $\det(D_W) = \prod \lambda_i = \prod |\lambda_i|^2 \ge 0$.
\end{enumerate}
This ensures that the probabilistic interpretation of the path integral is valid throughout the interpolation.
\end{proof}

\section{Uniform Bound on the Free Energy}
Using the positivity, we can establish the Vafa-Witten bound on the free energy density.
\begin{theorem}[Vafa-Witten Bound]
The vacuum energy density $E_0(M)$ of the vector-like theory is minimized at zero $\theta$-angle and provides a uniform lower bound for the spectral gap in the sense that no phase transition associated with parity breaking occurs.
\end{theorem}
This result rules out the scenario where the gap collapses due to the spontaneous breaking of vector symmetries (like Parity) along the interpolation path.
